% --------------------------------------------------
% Chapter 08
% Mini Project
% Credit Scoring MLOps Platform
% --------------------------------------------------

\label{chap:mini-project}

\section{Learning Objectives}

After completing this project, learners will be able to:

\begin{itemize}
    \item Apply the concepts introduced throughout Module 1.
    \item Design a complete MLOps architecture.
    \item Identify the major components of a production ML platform.
    \item Define roles and responsibilities.
    \item Establish governance and monitoring controls.
\end{itemize}

\section{Project Overview}

This project serves as the capstone exercise for Module 1.

Participants must design an enterprise-grade MLOps platform supporting a credit scoring model used by a retail bank.

The project integrates:

\begin{itemize}
    \item Machine Learning Lifecycle concepts
    \item MLOps principles
    \item Architecture design
    \item Governance requirements
    \item Operational considerations
\end{itemize}

The resulting platform will serve as the reference use case for the remaining modules of the training program.

\section{Business Context}

A retail bank receives thousands of credit applications each month.

The bank currently relies on manual decision-making processes and legacy scoring rules.

Management wants to deploy a machine learning solution capable of:

\begin{itemize}
    \item Improving approval decisions
    \item Reducing default rates
    \item Accelerating response times
    \item Increasing operational efficiency
\end{itemize}

The solution must satisfy both business and regulatory requirements.

\section{Business Requirements}

The machine learning platform must support:

\begin{itemize}
    \item Credit application scoring
    \item Model retraining
    \item Performance monitoring
    \item Regulatory audits
    \item Explainability requirements
\end{itemize}

Key stakeholders include:

\begin{itemize}
    \item Credit Risk Department
    \item Retail Banking
    \item Compliance
    \item Internal Audit
    \item IT Operations
\end{itemize}

\section{Available Data Sources}

The bank provides access to several data sources.

\subsection{Customer Data}

Examples include:

\begin{itemize}
    \item Age
    \item Income
    \item Employment status
    \item Marital status
\end{itemize}

\subsection{Loan Data}

Examples include:

\begin{itemize}
    \item Requested amount
    \item Loan purpose
    \item Loan duration
    \item Interest rate
\end{itemize}

\subsection{Credit Bureau Data}

Examples include:

\begin{itemize}
    \item Existing debt
    \item Payment history
    \item Credit utilization
    \item Number of inquiries
\end{itemize}

\subsection{Macroeconomic Data}

Examples include:

\begin{itemize}
    \item Unemployment rate
    \item Inflation
    \item Interest rate environment
\end{itemize}

\section{Target Architecture}

Participants must design a complete MLOps platform.

The reference architecture includes the following layers:

\begin{center}
Data Sources
\end{center}

\begin{center}
$\downarrow$
\end{center}

\begin{center}
Data Pipelines
\end{center}

\begin{center}
$\downarrow$
\end{center}

\begin{center}
Feature Store
\end{center}

\begin{center}
$\downarrow$
\end{center}

\begin{center}
Training Pipeline
\end{center}

\begin{center}
$\downarrow$
\end{center}

\begin{center}
Experiment Tracking
\end{center}

\begin{center}
$\downarrow$
\end{center}

\begin{center}
Model Registry
\end{center}

\begin{center}
$\downarrow$
\end{center}

\begin{center}
Deployment Layer
\end{center}

\begin{center}
$\downarrow$
\end{center}

\begin{center}
Monitoring Layer
\end{center}

\begin{center}
$\downarrow$
\end{center}

\begin{center}
Governance and Responsible AI
\end{center}

\section{Data Layer}

\subsection{Objectives}

The Data Layer must provide:

\begin{itemize}
    \item Reliable data ingestion
    \item Data quality controls
    \item Data lineage
    \item Dataset versioning
\end{itemize}

\subsection{Proposed Technologies}

Examples include:

\begin{itemize}
    \item Azure Data Factory
    \item Databricks
    \item Delta Lake
    \item DVC
\end{itemize}

\subsection{Deliverables}

Participants should define:

\begin{itemize}
    \item Data sources
    \item Data ownership
    \item Validation rules
    \item Refresh schedules
\end{itemize}

\section{Feature Engineering Layer}

\subsection{Objectives}

The Feature Layer must ensure:

\begin{itemize}
    \item Feature consistency
    \item Feature reuse
    \item Version management
    \item Metadata documentation
\end{itemize}

\subsection{Example Features}

Potential credit scoring features include:

\begin{itemize}
    \item Debt-to-income ratio
    \item Credit utilization
    \item Number of active loans
    \item Delinquency count
    \item Average account age
\end{itemize}

\subsection{Feature Store Requirements}

The Feature Store should support:

\begin{itemize}
    \item Offline training features
    \item Online serving features
    \item Metadata catalog
    \item Version control
\end{itemize}

\section{Training Environment}

\subsection{Objectives}

The training environment should support:

\begin{itemize}
    \item Reproducible experiments
    \item Hyperparameter tuning
    \item Automated training
    \item Performance evaluation
\end{itemize}

\subsection{Candidate Algorithms}

Examples include:

\begin{itemize}
    \item Logistic Regression
    \item Decision Trees
    \item Random Forest
    \item XGBoost
    \item LightGBM
\end{itemize}

\subsection{Evaluation Metrics}

Participants should evaluate:

\begin{itemize}
    \item Accuracy
    \item Precision
    \item Recall
    \item ROC-AUC
    \item KS Statistic
\end{itemize}

\section{Experiment Tracking}

The platform must record:

\begin{itemize}
    \item Parameters
    \item Metrics
    \item Artifacts
    \item Dataset versions
    \item Feature versions
\end{itemize}

\subsection{Recommended Technology}

\begin{itemize}
    \item MLflow Tracking
\end{itemize}

\section{Model Registry}

Participants must define:

\begin{itemize}
    \item Versioning strategy
    \item Approval workflow
    \item Promotion process
    \item Retirement policy
\end{itemize}

\subsection{Recommended Technology}

\begin{itemize}
    \item MLflow Model Registry
    \item Azure ML Registry
\end{itemize}

\section{Deployment Strategy}

\subsection{Objectives}

The deployment layer must provide:

\begin{itemize}
    \item Reliable inference services
    \item Controlled releases
    \item Scalability
    \item Security
\end{itemize}

\subsection{Deployment Patterns}

Participants should compare:

\begin{itemize}
    \item Batch Scoring
    \item Real-Time APIs
    \item Streaming Inference
\end{itemize}

\subsection{Recommended Architecture}

\begin{center}
Client Application
\end{center}

\begin{center}
$\downarrow$
\end{center}

\begin{center}
REST API
\end{center}

\begin{center}
$\downarrow$
\end{center}

\begin{center}
Credit Scoring Service
\end{center}

\begin{center}
$\downarrow$
\end{center}

\begin{center}
Prediction Response
\end{center}

\subsection{Containerization}

The scoring service should be packaged using:

\begin{itemize}
    \item Docker
\end{itemize}

Benefits include:

\begin{itemize}
    \item Portability
    \item Reproducibility
    \item Environment consistency
\end{itemize}

\subsection{Serving Framework}

Recommended technologies:

\begin{itemize}
    \item FastAPI
    \item Azure ML Managed Endpoints
\end{itemize}

\section{CI/CD Pipeline}

\subsection{Objectives}

The platform should automate:

\begin{itemize}
    \item Testing
    \item Validation
    \item Packaging
    \item Deployment
\end{itemize}

\subsection{Reference Workflow}

\begin{center}
Git Commit
\end{center}

\begin{center}
$\downarrow$
\end{center}

\begin{center}
Automated Tests
\end{center}

\begin{center}
$\downarrow$
\end{center}

\begin{center}
Model Validation
\end{center}

\begin{center}
$\downarrow$
\end{center}

\begin{center}
Container Build
\end{center}

\begin{center}
$\downarrow$
\end{center}

\begin{center}
Deployment
\end{center}

\subsection{Recommended Technologies}

\begin{itemize}
    \item GitHub Actions
    \item Azure DevOps
\end{itemize}

\section{Monitoring Framework}

\subsection{Objectives}

The monitoring framework must detect:

\begin{itemize}
    \item Performance degradation
    \item Data drift
    \item Concept drift
    \item Operational incidents
\end{itemize}

\subsection{Operational Monitoring}

Metrics include:

\begin{itemize}
    \item API latency
    \item Availability
    \item Error rates
    \item Resource consumption
\end{itemize}

\subsection{Model Monitoring}

Metrics include:

\begin{itemize}
    \item ROC-AUC
    \item Precision
    \item Recall
    \item Calibration
\end{itemize}

\subsection{Data Monitoring}

Examples include:

\begin{itemize}
    \item Missing value rates
    \item Distribution shifts
    \item Schema violations
\end{itemize}

\subsection{Recommended Technologies}

\begin{itemize}
    \item Azure Monitor
    \item Grafana
    \item Prometheus
    \item Evidently AI
\end{itemize}

\section{Responsible AI Requirements}

\subsection{Objectives}

The platform must support trustworthy AI practices.

Key objectives include:

\begin{itemize}
    \item Fairness
    \item Explainability
    \item Transparency
    \item Accountability
\end{itemize}

\subsection{Fairness Assessment}

Participants should evaluate:

\begin{itemize}
    \item Gender bias
    \item Age bias
    \item Group disparities
\end{itemize}

\subsection{Explainability}

The platform should provide:

\begin{itemize}
    \item Global explanations
    \item Local explanations
    \item Feature importance reports
\end{itemize}

\subsection{Recommended Technologies}

\begin{itemize}
    \item SHAP
    \item Fairlearn
    \item Responsible AI Dashboard
\end{itemize}

\section{Governance Framework}

\subsection{Objectives}

The platform must satisfy internal and external governance requirements.

Governance controls should ensure that models remain:

\begin{itemize}
    \item Traceable
    \item Auditable
    \item Explainable
    \item Compliant
\end{itemize}

\subsection{Model Documentation}

Each model should be accompanied by documentation describing:

\begin{itemize}
    \item Business purpose
    \item Data sources
    \item Feature definitions
    \item Training methodology
    \item Validation results
    \item Known limitations
\end{itemize}

\subsection{Approval Process}

A formal approval process should be established before production deployment.

Typical stakeholders include:

\begin{itemize}
    \item Business Owner
    \item Data Scientist
    \item Model Validator
    \item Risk Manager
    \item ML Engineer
\end{itemize}

\subsection{Audit Requirements}

The platform should retain:

\begin{itemize}
    \item Dataset versions
    \item Model versions
    \item Validation reports
    \item Deployment history
    \item Monitoring records
\end{itemize}

\subsection{Regulatory Considerations}

Examples include:

\begin{itemize}
    \item SR 11-7
    \item ECB TRIM
    \item EBA Guidelines
    \item EU AI Act
\end{itemize}

\section{Project Deliverables}

Participants should produce the following deliverables.

\subsection{Architecture Documentation}

The architecture document should describe:

\begin{itemize}
    \item Platform components
    \item Data flows
    \item Deployment strategy
    \item Monitoring controls
    \item Governance framework
\end{itemize}

\subsection{Technology Selection Report}

The report should justify:

\begin{itemize}
    \item Chosen tools
    \item Alternative solutions
    \item Architectural trade-offs
\end{itemize}

\subsection{Governance Package}

The governance package should contain:

\begin{itemize}
    \item Model inventory
    \item Validation framework
    \item Approval workflow
    \item Monitoring policy
\end{itemize}

\section{Evaluation Criteria}

Projects may be evaluated according to the following criteria.

\begin{table}[h]
\centering
\begin{tabular}{p{8cm}p{2cm}}
\toprule
Criterion & Weight \\
\midrule
Architecture Quality & 25\% \\
MLOps Principles & 20\% \\
Monitoring Design & 15\% \\
Governance Framework & 15\% \\
Responsible AI Controls & 15\% \\
Documentation Quality & 10\% \\
\bottomrule
\end{tabular}
\caption{Illustrative evaluation criteria}
\end{table}

\section{Final Reference Architecture}

The final architecture may be summarized as:

\begin{center}
Data Sources
\end{center}

\begin{center}
$\downarrow$
\end{center}

\begin{center}
Data Pipelines
\end{center}

\begin{center}
$\downarrow$
\end{center}

\begin{center}
Feature Store
\end{center}

\begin{center}
$\downarrow$
\end{center}

\begin{center}
Training Pipeline
\end{center}

\begin{center}
$\downarrow$
\end{center}

\begin{center}
MLflow Tracking
\end{center}

\begin{center}
$\downarrow$
\end{center}

\begin{center}
Model Registry
\end{center}

\begin{center}
$\downarrow$
\end{center}

\begin{center}
Docker Container
\end{center}

\begin{center}
$\downarrow$
\end{center}

\begin{center}
CI/CD Pipeline
\end{center}

\begin{center}
$\downarrow$
\end{center}

\begin{center}
Azure ML Endpoint
\end{center}

\begin{center}
$\downarrow$
\end{center}

\begin{center}
Monitoring
\end{center}

\begin{center}
$\downarrow$
\end{center}

\begin{center}
Responsible AI Controls
\end{center}

\begin{center}
$\downarrow$
\end{center}

\begin{center}
Governance Framework
\end{center}

\section{Summary}

This project integrates all concepts introduced in Module 1.

Participants have designed a complete machine learning platform covering:

\begin{itemize}
    \item Data Management
    \item Feature Engineering
    \item Model Development
    \item Experiment Tracking
    \item Model Registry
    \item Deployment
    \item Monitoring
    \item Responsible AI
    \item Governance
\end{itemize}

This architecture serves as the foundation for the remainder of the training program.

In subsequent modules, each component will be implemented in detail using technologies such as:

\begin{itemize}
    \item GitHub
    \item DVC
    \item Docker
    \item MLflow
    \item Azure Machine Learning
    \item GitHub Actions
    \item Prometheus
    \item Grafana
\end{itemize}

The learner now possesses a complete conceptual understanding of enterprise MLOps and is ready to move from architecture design to implementation.
